% ═══════════════════════════════════════════════════════════════════════════
%  فصل ۷ — ایران معاصر: مشروطه، پهلوی، جمهوری اسلامی
% ═══════════════════════════════════════════════════════════════════════════
\chapter{ایران معاصر: از رضاشاه تا جمهوری اسلامی}
\label{ch:modern-iran}

\begin{tcolorbox}[mybox, title={چکیده‌ی فصل}]
این فصل به بررسی مداخلات خارجی و نجات‌طلبی در ایران معاصر (۱۳۰۴ ش / ۱۹۲۵ م تا کنون) می‌پردازد. محورهای اصلی عبارت‌اند از: (۱) اشغال ایران در جنگ جهانی دوم و بحران آذربایجان؛ (۲) کودتای ۲۸ مرداد ۱۳۳۲ و نقش سیا/ام‌آی‌سیکس؛ (۳) انقلاب ۱۳۵۷ و مسئله‌ی «صدور انقلاب» به‌عنوان نجات‌طلبی معکوس؛ (۴) تحولات پس از ۱۳۵۷ و ادعاهای مداخله‌ی خارجی. تحلیل بر مبنای اسناد آرشیوی منتشرشده (سیا، \lr{MI6}، آرشیو ملی بریتانیا) و پژوهش‌های آکادمیک صورت می‌گیرد.
\end{tcolorbox}

% ─────────────────────────────────────────────────────────────────────────
\section{رضاشاه و تمرکزگرایی: پایان نجات‌طلبی‌های محلی؟}
\label{sec:ch07-reza-shah}
% ─────────────────────────────────────────────────────────────────────────

\subsection{سرکوب خوانین و یکپارچه‌سازی}
\label{subsec:ch07-centralization}

\histref{رضاشاه پهلوی}{۱۳۰۴–۱۳۲۰ ش / ۱۹۲۵–۱۹۴۱ م} با سیاست تمرکزگرایی شدید، تقریباً تمامی خوانین محلی، سران عشایر و حکام نیمه‌مستقل را سرکوب یا مطیع کرد. این شامل:
\begin{itemize}[nosep, rightmargin=1em]
  \item \textbf{شیخ خزعل} (خوزستان، ۱۳۰۳ ش): سرکوب نظامی.
  \item \textbf{سمکو آقا} (کردستان، ۱۳۰۱–۱۳۰۹ ش): سرکوب و ترور.%
  \footnote{\lr{Cronin, Stephanie. \textit{Tribal Politics in Iran}. Routledge, 2007, pp.~45--78.}}
  \item \textbf{خوانین بختیاری} (لرستان و چهارمحال): خلع سلاح و تبعید.
  \item \textbf{خوانین قشقایی} (فارس): مهار نسبی.%
  \footnote{بیات، کاوه. \textit{شورش عشایری فارس}. شیرازه، ۱۳۸۶.}
\end{itemize}

\begin{tcolorbox}[defbox, title={پارادوکس تمرکزگرایی رضاشاه}]
تمرکزگرایی رضاشاه از یک‌سو نجات‌طلبی‌های محلی را سرکوب کرد (شیخ خزعل دیگر نمی‌توانست از بریتانیا یاری بخواهد)، اما از سوی دیگر، با ایجاد نارضایتی‌های جدید (سرکوب فرهنگی اقوام، استبداد سیاسی، ناسیونالیسم تحمیلی)، زمینه‌ی نجات‌طلبی‌های آینده را فراهم کرد — به‌ویژه در آذربایجان و کردستان.

این همان الگوی «حذف واسطه‌ها با پیامدهای ناخواسته» است که در مورد خسرو دوم و لخمیان نیز مشاهده شد (فصل ۴).
\end{tcolorbox}

% ─────────────────────────────────────────────────────────────────────────
\section{اشغال ایران در جنگ جهانی دوم (۱۳۲۰–۱۳۲۵ ش)}
\label{sec:ch07-wwii}
% ─────────────────────────────────────────────────────────────────────────

\subsection{اشغال بدون دعوت}
\label{subsec:ch07-occupation}

اشغال ایران توسط متفقین (شهریور ۱۳۲۰ / سپتامبر ۱۹۴۱) نمونه‌ی بارز \concept{مداخله‌ی یک‌جانبه بدون دعوت} بود. هیچ‌بخشی از جامعه‌ی ایرانی — نه حکومت و نه اپوزیسیون — از متفقین درخواست اشغال نکرده بود.%
\footnote{\lr{Fawcett, Louise L'Estrange. \textit{Iran and the Cold War: The Azerbaijan Crisis of 1946}. Cambridge UP, 1992, pp.~18--35.}}

بهانه‌ی رسمی: جلوگیری از نفوذ آلمان نازی. اما نتیجه‌ی عملی:
\begin{itemize}[nosep, rightmargin=1em]
  \item برکناری رضاشاه و تبعید او.
  \item تقسیم ایران به مناطق اشغالی شوروی (شمال) و بریتانیا (جنوب).
  \item استفاده از ایران به‌عنوان «پل پیروزی» برای ارسال تجهیزات به شوروی.
  \item قحطی و تورم شدید.%
  \footnote{صدر هاشمی، محمد. \textit{تاریخ جراید و مجلات ایران}. ج ۴. کمال، ۱۳۶۳.}
\end{itemize}

\subsection{بحران آذربایجان و جمهوری خودمختار (۱۳۲۴–۱۳۲۵ ش)}
\label{subsec:ch07-azerbaijan-crisis}

\begin{tcolorbox}[casebox, title={بحران آذربایجان: نجات‌طلبی یا دست‌نشاندگی؟}]
در آذر ۱۳۲۴ ش / دسامبر ۱۹۴۵ م، با حمایت مستقیم ارتش سرخ شوروی، \histref{فرقه‌ی دموکرات آذربایجان}{به رهبری سیدجعفر پیشه‌وری} «حکومت ملی آذربایجان» را در تبریز اعلام کرد. هم‌زمان، در کردستان نیز \histref{جمهوری مهاباد}{به رهبری قاضی محمد} تأسیس شد.%
\footnote{\lr{Fawcett, 1992, pp.~89--120.}}

\textbf{پرسش کلیدی:} آیا این جنبش‌ها «نجات‌طلبی» واقعی (بیان خواست مردم آذربایجان/کردستان) بودند یا «دست‌نشاندگی» شوروی؟

\textbf{شواهد نجات‌طلبی واقعی:}
\begin{itemize}[nosep, rightmargin=1em]
  \item نارضایتی واقعی از تبعیض زبانی و فرهنگی دوره‌ی رضاشاه.
  \item خواست خودمختاری (نه لزوماً جدایی) بخشی از نخبگان آذری.%
  \footnote{\lr{Atabaki, Touraj. \textit{Azerbaijan: Ethnicity and the Struggle for Power in Iran}. I.B. Tauris, 2000, pp.~85--115.}}
  \item اصلاحات اجتماعی فرقه (اصلاحات ارضی، آموزش به زبان ترکی) مورد استقبال بخشی از مردم بود.
\end{itemize}

\textbf{شواهد دست‌نشاندگی:}
\begin{itemize}[nosep, rightmargin=1em]
  \item فرقه بدون حمایت ارتش سرخ توانایی تأسیس نداشت.
  \item شوروی از فرقه به‌عنوان ابزار فشار برای گرفتن امتیاز نفت شمال استفاده کرد.%
  \footnote{\lr{Hasanli, Jamil. \textit{At the Dawn of the Cold War: The Soviet-American Crisis over Iranian Azerbaijan, 1941--1946}. Rowman and Littlefield, 2006.}}
  \item پس از عقب‌نشینی شوروی (آذر ۱۳۲۵ ش)، فرقه بدون مقاومت جدی فروپاشید — نشان‌دهنده‌ی فقدان پایگاه اجتماعی عمیق.
  \item پیشه‌وری خود سابقه‌ی طولانی در حزب کمونیست شوروی داشت.%
  \footnote{آبراهامیان، یرواند. \textit{ایران بین دو انقلاب}. ترجمه‌ی احمد گل‌محمدی و محمدابراهیم فتاحی. نی، ۱۳۷۷، ص ۲۵۵--۲۸۰.}
\end{itemize}

\textbf{نتیجه‌ی تحلیل:} بحران آذربایجان نمونه‌ی \concept{نجات‌طلبی تحریک‌شده} (\lr{Induced Salvation-Seeking}) بود: نارضایتی واقعی وجود داشت اما سازمان‌دهی و تقویت آن مستقیماً توسط قدرت خارجی (شوروی) صورت گرفت. هنگامی که قدرت خارجی عقب کشید، جنبش فروپاشید — تأیید قاعده‌ی «بی‌وفایی حامی خارجی».
\end{tcolorbox}

\begin{figure}[htbp]
\centering
\begin{tikzpicture}[
  every node/.style={font=\scriptsize},
  x=0.018\linewidth,
]

% خط زمان
\draw[line width=2pt, PrimaryMid] (0,0) -- (120,0);

% رویدادها
\node[circle, fill=AccentRed, minimum size=7pt, inner sep=0pt] at (0,0) {};
\node[below=0.5cm, text width=2.5cm, align=center] at (0,0)
  {\rl{شهریور ۱۳۲۰}\\%
   \rl{اشغال ایران}};

\node[circle, fill=AccentAmber, minimum size=7pt, inner sep=0pt] at (30,0) {};
\node[above=0.5cm, text width=2.5cm, align=center] at (30,0)
  {\rl{شهریور ۱۳۲۲}\\%
   \rl{اعلامیه‌ی تهران}};

\node[circle, fill=EraContemp, minimum size=7pt, inner sep=0pt] at (55,0) {};
\node[below=0.5cm, text width=2.5cm, align=center] at (55,0)
  {\rl{آذر ۱۳۲۴}\\%
   \rl{تأسیس فرقه}};

\node[circle, fill=AccentGreen, minimum size=7pt, inner sep=0pt] at (75,0) {};
\node[above=0.5cm, text width=2.5cm, align=center] at (75,0)
  {\rl{اردیبهشت ۱۳۲۵}\\%
   \rl{توافق قوام-شوروی}};

\node[circle, fill=AccentRed, minimum size=7pt, inner sep=0pt] at (95,0) {};
\node[below=0.5cm, text width=2.5cm, align=center] at (95,0)
  {\rl{آذر ۱۳۲۵}\\%
   \rl{سقوط فرقه}};

\node[circle, fill=NeutralDark, minimum size=7pt, inner sep=0pt] at (115,0) {};
\node[above=0.5cm, text width=2.5cm, align=center] at (115,0)
  {\rl{مهر ۱۳۲۶}\\%
   \rl{رد امتیاز نفت شمال}};

\end{tikzpicture}
\caption{خط زمانی بحران آذربایجان (۱۳۲۰–۱۳۲۶ ش)}
\label{fig:ch07-azerbaijan-timeline}
\end{figure}

% ─────────────────────────────────────────────────────────────────────────
\section{کودتای ۲۸ مرداد ۱۳۳۲: مداخله‌ی خارجی بدون نجات‌طلبی مردمی}
\label{sec:ch07-coup}
% ─────────────────────────────────────────────────────────────────────────

\subsection{زمینه‌ها و اسناد}
\label{subsec:ch07-coup-background}

کودتای ۲۸ مرداد ۱۳۳۲ / ۱۹ اوت ۱۹۵۳ علیه دولت \histref{دکتر محمد مصدق}{نخست‌وزیر ایران} یکی از مستندترین نمونه‌های مداخله‌ی خارجی در تاریخ معاصر است. اسناد آزادشده‌ی سیا (۲۰۱۳ و ۲۰۱۷) و آرشیو ملی بریتانیا، جزئیات عملیات \lr{TPAJAX} (سیا) و \lr{Operation Boot} (\lr{MI6}) را فاش کرده‌اند.%
\footnote{\lr{Gasiorowski, Mark J., and Malcolm Byrne, eds. \textit{Mohammad Mosaddeq and the 1953 Coup in Iran}. Syracuse UP, 2004.}}

\begin{legalNote}{اسناد سیا درباره‌ی کودتای ۱۹۵۳}
گزارش داخلی سیا (تاریخ‌نگاری عملیات \lr{TPAJAX}، تنظیم: دونالد ویلبر، ۱۹۵۴) تصریح می‌کند:

«مشخص بود که بدون حمایت بریتانیا و ایالات متحده، مخالفان مصدق — نه شاه و نه ارتش و نه روحانیون — توانایی سرنگونی او را نداشتند ... عملیات شامل رشوه‌دادن به روزنامه‌نگاران، روحانیون، افسران ارتش و لمپن‌های خیابانی بود.»%
\footnote{\lr{CIA. "Clandestine Service History: Overthrow of Premier Mossadeq of Iran." Written by Donald N. Wilber, March 1954. Declassified 2017. Available at National Security Archive, George Washington University.}}
\end{legalNote}

\subsection{تحلیل شش‌وجهی}
\label{subsec:ch07-coup-analysis}

\begin{tcolorbox}[casebox, title={تحلیل شش‌وجهی: کودتای ۲۸ مرداد ۱۳۳۲}]
\begin{itemize}[nosep, rightmargin=1em]
  \item \textbf{ساختاری:} مصدق نخست‌وزیر منتخب مجلس بود و حمایت مردمی گسترده‌ای داشت. اما اختلاف با شاه، بخشی از ارتش و برخی روحانیون (آیت‌الله کاشانی) وجود داشت.

  \item \textbf{حقوقی:} مداخله از هیچ منظر حقوقی بین‌المللی مشروعیت نداشت. نه مجوز شورای امنیت وجود داشت و نه «دعوت» از حکومت مشروع (مصدق خود حکومت مشروع بود).

  \item \textbf{اقتصادی:} ملی‌شدن نفت (۱۳۲۹ ش) محرک اصلی بریتانیا بود. شرکت نفت ایران و انگلیس (\lr{AIOC}) سالانه سودی بیش از کل بودجه‌ی دولت ایران داشت.%
  \footnote{\lr{Elm, Mostafa. \textit{Oil, Power, and Principle: Iran's Oil Nationalization and Its Aftermath}. Syracuse UP, 1992.}}

  \item \textbf{نظامی:} کودتا عمدتاً با ابزارهای اطلاعاتی (رشوه، تبلیغات، ایجاد آشوب) انجام شد. مداخله‌ی نظامی مستقیم لازم نبود.

  \item \textbf{روان‌شناختی:} بخشی از نخبگان ایرانی (محافظه‌کاران سلطنت‌طلب، برخی روحانیون و نظامیان) «نجات‌طلبی معکوس» داشتند — یعنی از قدرت خارجی خواستند تا مصدق را سرنگون کند. اما این «نجات‌طلبی» از سوی توده‌ی مردم نبود.%
  \footnote{\lr{Abrahamian, Ervand. \textit{The Coup: 1953, the CIA, and the Roots of Modern U.S.-Iranian Relations}. New Press, 2013.}}

  \item \textbf{ژئوپولیتیک:} جنگ سرد، نفت و ترس از «سقوط ایران به دامان کمونیسم» (دکترین دومینو) محرک‌های ایالات متحده بود.

  \item \textbf{نتیجه (طیف):} \textbf{D–E}. بازگشت شاه و تحکیم استبداد سلطنتی؛ وابستگی ساختاری به غرب؛ سرکوب جنبش ملی؛ بذر انقلاب ۱۳۵۷.
\end{itemize}
\end{tcolorbox}

\subsection{کودتای ۲۸ مرداد از منظر «نجات‌طلبی»}
\label{subsec:ch07-coup-salvation}

\begin{table}[htbp]
\centering
\caption{بازیگران داخلی کودتا و رابطه‌شان با قدرت خارجی}\label{tab:ch07-coup-actors}
\begin{adjustbox}{max width=\linewidth}
\begin{tabular}{
  >{\raggedleft\arraybackslash\bfseries}p{2.5cm}
  >{\raggedleft\arraybackslash}p{3cm}
  >{\raggedleft\arraybackslash}p{3cm}
  >{\raggedleft\arraybackslash}p{2.5cm}
  >{\raggedleft\arraybackslash}p{2.5cm}}
\toprule
\rowcolor{PrimaryDeep}
\textcolor{white}{بازیگر} &
\textcolor{white}{نوع رابطه با خارجی} &
\textcolor{white}{انگیزه} &
\textcolor{white}{آگاهی از پیامد} &
\textcolor{white}{سرنوشت} \\
\midrule

محمدرضاشاه &
همکاری فعال با سیا/\lr{MI6} &
حفظ سلطنت &
ناآگاه از عمق وابستگی &
۲۵ سال حکومت؛ انقلاب ۵۷%
\footnote{\lr{Milani, Abbas. \textit{The Shah}. Palgrave Macmillan, 2011.}} \\
\rowcolor{NeutralLight}
سرلشکر زاهدی &
مجری کودتا؛ تماس مستقیم با سیا &
قدرت‌طلبی شخصی &
نسبتاً آگاه &
نخست‌وزیری و سپس تبعید \\
آیت‌الله کاشانی &
همسویی موقت (نه لزوماً تماس مستقیم) &
اختلاف شخصی با مصدق &
احتمالاً ناآگاه &
حاشیه‌نشینی سیاسی%
\footnote{آبراهامیان، ۱۳۷۷، ص ۳۵۰--۳۷۰.} \\
\rowcolor{NeutralLight}
شعبان بی‌مخ و لمپن‌ها &
ابزار سیا (پرداخت مستقیم) &
مالی &
کاملاً ناآگاه &
حاشیه \\
تودهای‌ها &
(متهم به نجات‌طلبی شوروی) &
ایدئولوژیک &
بحث‌برانگیز &
سرکوب شدید \\

\bottomrule
\end{tabular}
\end{adjustbox}
\end{table}

\begin{tcolorbox}[defbox, title={نکته‌ی کلیدی: تمایز «همکاری نخبگان» و «نجات‌طلبی مردمی»}]
کودتای ۲۸ مرداد نمونه‌ای است از \concept{همکاری نخبگان با مداخله‌گر خارجی} — نه نجات‌طلبی مردمی. توده‌ی مردم نه خواهان کودتا بودند و نه از آن اطلاع داشتند. تمایز «همکاری نخبگان» از «نجات‌طلبی مردمی» اهمیت تحلیلی فراوانی دارد:

\begin{itemize}[nosep, rightmargin=1em]
  \item \textbf{همکاری نخبگان:} تعداد محدودی از صاحبان قدرت (نظامی، مذهبی، سلطنتی) با مداخله‌گر خارجی علیه حکومت منتخب همکاری می‌کنند. پایگاه اجتماعی ضعیف.
  \item \textbf{نجات‌طلبی مردمی:} بخش قابل‌توجهی از جمعیت خواهان مداخله‌ی خارجی برای رهایی از ستم‌اند. پایگاه اجتماعی قوی.
\end{itemize}
کودتای ۲۸ مرداد قطعاً در دسته‌ی اول قرار دارد.
\end{tcolorbox}

% ─────────────────────────────────────────────────────────────────────────
\section{انقلاب ۱۳۵۷ و «صدور انقلاب»: نجات‌طلبی معکوس}
\label{sec:ch07-revolution}
% ─────────────────────────────────────────────────────────────────────────

\subsection{انقلاب ۱۳۵۷ از منظر مداخله‌ی خارجی}
\label{subsec:ch07-revolution-foreign}

انقلاب ۱۳۵۷ یکی از نادرترین انقلاب‌های بزرگ تاریخ معاصر است که \textbf{بدون کمک نظامی خارجی مستقیم} و عمدتاً با بسیج مردمی صورت گرفت. با این حال، نقش عوامل خارجی محل بحث است:

\begin{longtable}{
  >{\raggedleft\arraybackslash\bfseries}p{2.5cm}
  >{\raggedleft\arraybackslash}p{4cm}
  >{\raggedleft\arraybackslash}p{3cm}
  >{\raggedleft\arraybackslash}p{3.5cm}}
\caption{نظریه‌های نقش عوامل خارجی در انقلاب ۱۳۵۷}\label{tab:ch07-revolution-theories}\\
\toprule
\rowcolor{PrimaryDeep}
\textcolor{white}{\textbf{نظریه}} &
\textcolor{white}{\textbf{شرح}} &
\textcolor{white}{\textbf{طرفداران}} &
\textcolor{white}{\textbf{شواهد}} \\
\midrule
\endfirsthead
\toprule
\rowcolor{PrimaryDeep}
\textcolor{white}{\textbf{نظریه}} &
\textcolor{white}{\textbf{شرح}} &
\textcolor{white}{\textbf{طرفداران}} &
\textcolor{white}{\textbf{شواهد}} \\
\midrule
\endhead
\bottomrule
\endlastfoot

توطئه‌ی غربی &
غرب (امریکا/بریتانیا) شاه را «رها» کرد تا ایران «سبز» (اسلام‌گرا) شود و سد نفوذ شوروی باشد &
سلطنت‌طلبان، برخی محافل ایرانی &
ضعیف؛ اسناد سیا نشان‌دهنده‌ی غافل‌گیری امریکاست%
\footnote{\lr{Sick, Gary. \textit{All Fall Down: America's Tragic Encounter with Iran}. Random House, 1985.}} \\
\rowcolor{NeutralLight}
توطئه‌ی شوروی &
شوروی و حزب توده انقلاب را هدایت کردند &
شاه (در خاطراتش) &
بسیار ضعیف؛ توده‌ای‌ها نقش حاشیه‌ای داشتند \\
انقلاب مردمی مستقل &
انقلاب نتیجه‌ی عوامل داخلی بود و مداخله‌ی خارجی نقش تعیین‌کننده‌ای نداشت &
اکثر پژوهشگران آکادمیک &
قوی؛ بسیج میلیونی مردم%
\footnote{\lr{Kurzman, Charles. \textit{The Unthinkable Revolution in Iran}. Harvard UP, 2004.}} \\
\rowcolor{NeutralLight}
عوامل خارجی تسهیل‌کننده &
غرب مانع سرکوب شد (فشار کارتر درباره‌ی حقوق بشر) — نه عامل انقلاب بلکه تسهیل‌کننده &
برخی محققان &
متوسط؛ سیاست حقوق بشر کارتر مستند است%
\footnote{\lr{Bill, James A. \textit{The Eagle and the Lion: The Tragedy of American-Iranian Relations}. Yale UP, 1988.}} \\

\end{longtable}

\subsection{صدور انقلاب: نجات‌طلبی معکوس}
\label{subsec:ch07-export-revolution}

پس از انقلاب ۱۳۵۷، جمهوری اسلامی سیاست \concept{صدور انقلاب} را در پیش گرفت — یعنی تلاش برای تکرار الگوی انقلاب اسلامی در سایر کشورها. از منظر این پژوهش، صدور انقلاب نوعی \concept{نجات‌طلبی معکوس} (\lr{Reverse Salvation-Seeking}) است: نه مردم تحت ستمی از ایران دعوت می‌کنند، بلکه ایران خود را «ناجی» ملت‌های مسلمان تحت ستم معرفی می‌کند.%
\footnote{\lr{Hunter, Shireen T. \textit{Iran and the World: Continuity in a Revolutionary Decade}. Indiana UP, 1990, pp.~35--55.}}

\begin{tcolorbox}[defbox, title={صدور انقلاب: مصادیق و نتایج}]
\begin{itemize}[nosep, rightmargin=1em]
  \item \textbf{لبنان:} تأسیس حزب‌الله (۱۹۸۲) با حمایت مستقیم سپاه پاسداران. موفق‌ترین نمونه‌ی صدور انقلاب.%
  \footnote{\lr{Norton, Augustus Richard. \textit{Hezbollah: A Short History}. Princeton UP, 2007.}}
  \item \textbf{عراق:} حمایت از احزاب شیعی (حزب الدعوه، مجلس اعلا). نتایج مختلط — هم نفوذ و هم مقاومت.
  \item \textbf{بحرین:} تلاش ناکام برای کودتا (۱۹۸۱). شکست.%
  \footnote{\lr{Chubin, Shahram, and Charles Tripp. \textit{Iran and Iraq at War}. I.B. Tauris, 1988.}}
  \item \textbf{افغانستان:} حمایت از هزاره‌ها و شیعیان. نتایج بلندمدت.
  \item \textbf{یمن:} حمایت از حوثی‌ها (پس از ۲۰۱۱). تداوم.
\end{itemize}
\end{tcolorbox}

% ─────────────────────────────────────────────────────────────────────────
\section{دوره‌ی پس از ۱۳۵۷: اسناد اطلاعاتی و مداخلات مستند}
\label{sec:ch07-post-revolution}
% ─────────────────────────────────────────────────────────────────────────

\subsection{عملیات‌های مستند سیا و \texorpdfstring{\lr{MI6}}{MI6} علیه جمهوری اسلامی}
\label{subsec:ch07-cia-mi6}

بر اساس اسناد آزادشده و گزارش‌های معتبر:

\begin{longtable}{
  >{\raggedleft\arraybackslash\bfseries}p{2.5cm}
  >{\raggedleft\arraybackslash}p{3cm}
  >{\raggedleft\arraybackslash}p{4cm}
  >{\raggedleft\arraybackslash}p{3.5cm}}
\caption{عملیات‌ها و طرح‌های مستند علیه جمهوری اسلامی}\label{tab:ch07-operations}\\
\toprule
\rowcolor{PrimaryDeep}
\textcolor{white}{\textbf{عملیات/طرح}} &
\textcolor{white}{\textbf{سال}} &
\textcolor{white}{\textbf{شرح}} &
\textcolor{white}{\textbf{منبع}} \\
\midrule
\endfirsthead
\toprule
\rowcolor{PrimaryDeep}
\textcolor{white}{\textbf{عملیات/طرح}} &
\textcolor{white}{\textbf{سال}} &
\textcolor{white}{\textbf{شرح}} &
\textcolor{white}{\textbf{منبع}} \\
\midrule
\endhead
\bottomrule
\endlastfoot

کودتای نوژه &
۱۳۵۹ ش &
تلاش ناکام افسران ارتش (با ادعای ارتباط با عراق و غرب) &
آبراهامیان، ۱۳۸۲%
\footnote{آبراهامیان، یرواند. \textit{تاریخ ایران مدرن}. ترجمه‌ی محمدابراهیم فتاحی. نی، ۱۳۸۹، ص ۲۱۵--۲۲۰.} \\
\rowcolor{NeutralLight}
حمایت از صدام &
۱۳۵۹–۱۳۶۷ ش &
کمک اطلاعاتی و تسلیحاتی غرب و عرب به عراق در جنگ هشت‌ساله &
\lr{Hiro, Dilip. \textit{The Longest War}. Routledge, 1991.}%
\footnote{\lr{Hiro, Dilip. \textit{The Longest War: The Iran-Iraq Military Conflict}. Routledge, 1991.}} \\
ایران‌گیت &
۱۳۶۵ ش &
فروش مخفیانه‌ی سلاح امریکایی به ایران (با وجود تحریم!) &
\lr{Walsh Report, 1993}%
\footnote{\lr{Walsh, Lawrence E. \textit{Iran-Contra: The Final Report}. Times Books, 1994.}} \\
\rowcolor{NeutralLight}
حمایت از گروه‌های اپوزیسیون &
۱۳۶۰–کنون &
حمایت‌های مختلف از سازمان مجاهدین، گروه‌های قومی و غیره &
\lr{Hersh, Seymour. "Preparing the Battlefield." \textit{New Yorker}, July 2008.}%
\footnote{\lr{Hersh, Seymour M. "Preparing the Battlefield." \textit{The New Yorker}, 7 July 2008.}} \\
\lr{Stuxnet} &
ح. ۱۳۸۸–۱۳۸۹ ش &
حمله‌ی سایبری به تأسیسات هسته‌ای نطنز &
\lr{Zetter, Kim. \textit{Countdown to Zero Day}. Crown, 2014.}%
\footnote{\lr{Zetter, Kim. \textit{Countdown to Zero Day: Stuxnet and the Launch of the World's First Digital Weapon}. Crown, 2014.}} \\

\end{longtable}

% ─────────────────────────────────────────────────────────────────────────
\section{نجات‌طلبی معاصر: گروه‌های قومی و مسئله‌ی «طلب مداخله»}
\label{sec:ch07-ethnic-salvation}
% ─────────────────────────────────────────────────────────────────────────

\subsection{کردها}
\label{subsec:ch07-kurds}

جنبش‌های کردی در ایران — از جمهوری مهاباد (۱۳۲۴ ش) تا حزب دموکرات کردستان ایران و کومه‌له — همواره با اتهام «نجات‌طلبی» مواجه بوده‌اند. واقعیت پیچیده‌تر است:

\begin{itemize}[nosep, rightmargin=1em]
  \item جمهوری مهاباد مستقیماً وابسته به شوروی بود (فصل قبل).
  \item پس از ۱۳۵۷، گروه‌های کردی اساساً خواهان \concept{خودمختاری} بودند نه جدایی — و از نیروی خارجی «دعوت» نکردند بلکه از حکومت مرکزی خواستار مذاکره شدند.%
  \footnote{\lr{Vali, Abbas. \textit{Kurds and the State in Iran}. I.B. Tauris, 2011.}}
  \item عراق (صدام) از گروه‌های کردی ایران علیه جمهوری اسلامی حمایت کرد — در همان حال که کردهای عراق را نسل‌کشی می‌کرد (انفال).
\end{itemize}

\subsection{بلوچ‌ها و عرب‌ها}
\label{subsec:ch07-baluch-arab}

گروه‌های مسلح بلوچ (\lr{Jundallah} / جندالله) و برخی گروه‌های جدایی‌طلب عرب خوزستانی در دوره‌های مختلف متهم به دریافت حمایت از عربستان سعودی، پاکستان یا امریکا بوده‌اند.%
\footnote{\lr{Hersh, 2008.}}

\begin{tcolorbox}[enemybox, title={نکته‌ی اخلاقی-تحلیلی}]
تفکیک \textbf{«حق اعتراض و خودمختاری‌خواهی»} از \textbf{«نجات‌طلبی از طریق قدرت خارجی»} ضروری است. بسیاری از گروه‌های قومی ایران خواست‌های مشروعی (حقوق زبانی، فرهنگی، عدم تبعیض) دارند که ربطی به مداخله‌ی خارجی ندارد. اتهام «وابستگی به بیگانه» اغلب ابزاری برای سرکوب هرگونه اعتراض مشروع بوده است.

با این حال، شواهد مستندی نیز وجود دارد که برخی قدرت‌های خارجی از نارضایتی‌های قومی ایران \textbf{بهره‌برداری ابزاری} کرده‌اند — بدون آنکه لزوماً به خواست‌های واقعی آن گروه‌ها پایبند باشند.
\end{tcolorbox}

% ─────────────────────────────────────────────────────────────────────────
\section{الگوهای مداخله‌ی خارجی در ایران معاصر: جدول تطبیقی}
\label{sec:ch07-patterns}
% ─────────────────────────────────────────────────────────────────────────

\begin{landscape}
\begin{table}[htbp]
\centering
\caption{جدول تطبیقی مداخلات خارجی در ایران معاصر}\label{tab:ch07-comparative}
\begin{adjustbox}{max width=\linewidth}
\begin{tabular}{
  >{\raggedleft\arraybackslash\bfseries\small}p{2.2cm}
  >{\raggedleft\arraybackslash\small}p{1.8cm}
  >{\raggedleft\arraybackslash\small}p{2cm}
  >{\raggedleft\arraybackslash\small}p{2.2cm}
  >{\raggedleft\arraybackslash\small}p{2cm}
  >{\raggedleft\arraybackslash\small}p{2cm}
  >{\raggedleft\arraybackslash\small}p{2cm}
  >{\raggedleft\arraybackslash\small}p{1.8cm}}
\toprule
\rowcolor{PrimaryDeep}
\textcolor{white}{مورد} &
\textcolor{white}{سال} &
\textcolor{white}{مداخله‌گر} &
\textcolor{white}{نوع مداخله} &
\textcolor{white}{دعوت داخلی؟} &
\textcolor{white}{نتیجه‌ی فوری} &
\textcolor{white}{نتیجه‌ی بلندمدت} &
\textcolor{white}{طیف} \\
\midrule

اشغال متفقین &
۱۳۲۰ &
شوروی/بریتانیا &
نظامی مستقیم &
خیر &
برکناری رضاشاه &
خروج (۱۳۲۵) اما آثار ماندگار &
C \\
\rowcolor{NeutralLight}
بحران آذربایجان &
۱۳۲۴–۲۵ &
شوروی &
نظامی + سیاسی &
نیمه (تحریک‌شده) &
حکومت خودمختار &
فروپاشی سریع &
E \\
کودتای ۲۸ مرداد &
۱۳۳۲ &
سیا/\lr{MI6} &
اطلاعاتی &
نخبگان (نه مردم) &
سقوط مصدق &
استبداد + انقلاب ۵۷ &
D–E \\
\rowcolor{NeutralLight}
حمایت از صدام &
۱۳۵۹–۶۷ &
غرب + عرب &
تسلیحاتی + اطلاعاتی &
خیر &
طولانی‌شدن جنگ &
یک میلیون کشته &
F \\
صدور انقلاب &
۱۳۵۸–کنون &
ایران &
ایدئولوژیک + نظامی &
بخشاً بله &
نتایج مختلط &
نفوذ منطقه‌ای &
C–D \\
\rowcolor{NeutralLight}
تحریم‌ها &
۱۳۵۸–کنون &
امریکا + اروپا &
اقتصادی &
خیر (اما برخی اپوزیسیون حمایت) &
فشار اقتصادی &
بحث‌برانگیز &
? \\

\bottomrule
\end{tabular}
\end{adjustbox}
\end{table}
\end{landscape}

% ─────────────────────────────────────────────────────────────────────────
\section{جمع‌بندی فصل}
\label{sec:ch07-summary}
% ─────────────────────────────────────────────────────────────────────────

\begin{tcolorbox}[wavebox, title={یافته‌های کلیدی فصل ۷}]
\begin{enumerate}[nosep, label=\textbf{\arabic*.}]
  \item \textbf{تمرکزگرایی رضاشاه} نجات‌طلبی‌های محلی را سرکوب کرد اما نارضایتی‌های جدید ایجاد نمود — الگوی «حذف واسطه‌ها».

  \item \textbf{بحران آذربایجان} نمونه‌ی \concept{نجات‌طلبی تحریک‌شده} بود: نارضایتی واقعی + سازمان‌دهی خارجی. فروپاشی پس از عقب‌نشینی شوروی نشان داد که بدون حامی خارجی، جنبش توان ادامه نداشت.

  \item \textbf{کودتای ۲۸ مرداد} نمونه‌ی \concept{همکاری نخبگان} با مداخله‌گر بود — نه نجات‌طلبی مردمی. پیامدش ۲۵ سال استبداد و سپس انقلاب بود.

  \item \textbf{انقلاب ۱۳۵۷} نمونه‌ی نادر انقلاب مستقل بدون مداخله‌ی نظامی خارجی بود — اما «صدور انقلاب» آن را به نوعی «نجات‌طلبی معکوس» تبدیل کرد.

  \item \textbf{دوره‌ی پس از انقلاب} ترکیبی از مداخلات خارجی (جنگ، تحریم، عملیات‌های مخفی) و ادعاهای نجات‌طلبی قومی را شامل می‌شود. تفکیک «اعتراض مشروع» از «نجات‌طلبی ابزاری» ضروری اما دشوار است.

  \item \textbf{الگوی ثابت:} در تمامی موارد بررسی‌شده، مداخله‌گر خارجی منافع خود را بر خواسته‌ی متحدان داخلی‌اش ترجیح داده است — از شوروی (آذربایجان) تا امریکا (شاه) تا بریتانیا (خزعل).
\end{enumerate}
\end{tcolorbox}

\cleardoublepage